% GENERATED FILE. Edit canonical lessons, then run npm run generate:books.
% canonical-source-hash: fnv1a64:e978d75d679aa31e
% canonical-lessons: SA-C259-asmakam, SA-C259-yusmakam, SA-C259-plural-possessors-practice

\chapter{Our and your-plural}
\label{ch:sa-our-and-your-plural}
\hlchaptermodality{259}

\begin{chapteropening}
\textbf{By the end of this chapter:} \emph{I can distinguish one possessor from several and say our and your-plural in Sanskrit.}

Complete chapter 259 by choosing and retrieving the singular and plural possessor forms in known noun frames.
\end{chapteropening}

\section[\texorpdfstring{\sk{अस्माकम्} (asmākam)}{अस्माकम् (asmākam)}]{\sk{अस्माकम्} (asmākam) — our, of us}

\label{lesson:SA-C259-asmakam}



\begin{quote}
Say \textbf{\sk{मम} \sk{गृहम्}} (\emph{mama gṛham}), “my house”. One person possesses the house. Keep \textbf{\sk{मम}} ready.
\end{quote}

\subsection*{You'll want to know: one possessor, then several}
\noindent
\begin{tabularx}{\linewidth}{@{}>{\raggedright\arraybackslash}X>{\raggedright\arraybackslash}X>{\raggedright\arraybackslash}X@{}}
\toprule
\textbf{possessors} & \textbf{Sanskrit} & \textbf{meaning} \\
\midrule
one: me & \textbf{\sk{मम}} (\emph{mama}) & my, of me \\
several: us & \textbf{\sk{अस्माकम्}} (\emph{asmākam}) & our, of us \\
\bottomrule
\end{tabularx}

Say the new form slowly: \emph{as-mā-kam}. The long \textbf{ā} holds the middle beat.

These are genitive forms: they answer “whose?” Sanskrit uses that relationship where English often uses a possessive word.

\begin{cousinweb}
Whitney's \emph{Sanskrit Grammar} §493 lists \textbf{\sk{अस्माकम्}} as the genitive plural in the first-person paradigm. Historically, he explains it as a fixed case-form of an older adjective stem: the form is derived from \emph{asmāka}. The history is unusual; the job in the sentence is simple: \textbf{\sk{अस्माकम्}} answers “of us / whose?” and stays fixed before the noun.

Source: William Dwight Whitney, \emph{Sanskrit Grammar} (1889), §§493 and 516.
\end{cousinweb}

\subsection*{Your turn}
Say these aloud:

\begin{itemize}
\raggedright
  \item “my house” — \emph{mama gṛham}
  \item “our house” — \emph{asmākam gṛham}
  \item one possessor — \emph{mama}; several possessors — \emph{asmākam}
\end{itemize}

\begin{tcolorbox}[breakable,colback=teal!4,colframe=teal!35!black,title={Before you move on}]
Several people point to one house and say “our house”. Which first word do they use? (\textbf{\sk{अस्माकम्}}, \emph{asmākam}.)
\end{tcolorbox}
\section[\texorpdfstring{\sk{युष्माकम्} (yuṣmākam)}{युष्माकम् (yuṣmākam)}]{\sk{युष्माकम्} (yuṣmākam) — your, of you all}

\label{lesson:SA-C259-yusmakam}



\begin{quote}
Say “our” (\textbf{\sk{अस्माकम्}}, \emph{asmākam}). Then recall the first word of \textbf{\sk{तव} \sk{नाम} \sk{किम्}?} (\emph{tava nāma kim?}): \textbf{\sk{तव}}, “your” to one familiar person.
\end{quote}

\subsection*{You'll want to know: one addressee, then several}
\noindent
\begin{tabularx}{\linewidth}{@{}>{\raggedright\arraybackslash}X>{\raggedright\arraybackslash}X>{\raggedright\arraybackslash}X@{}}
\toprule
\textbf{people addressed} & \textbf{Sanskrit} & \textbf{meaning} \\
\midrule
one & \textbf{\sk{तव}} (\emph{tava}) & your, of you \\
several & \textbf{\sk{युष्माकम्}} (\emph{yuṣmākam}) & your, of you all \\
\bottomrule
\end{tabularx}

Say the new form slowly: \emph{yuṣ-mā-kam}. The ending sounds like \emph{asmākam}; the opening tells you that the possessors are the people being addressed.

Like \emph{tava}, \textbf{\sk{युष्माकम्}} is a genitive. It answers “whose?” for several people, and the noun it comes before does not have to change.

\begin{cousinweb}
Whitney's \emph{Sanskrit Grammar} §493 places \textbf{\sk{युष्माकम्}} beside \textbf{\sk{अस्माकम्}} as the second-person genitive plural. It too is historically a fixed case-form of an older adjective stem: the form is derived from \emph{yuṣmāka}. For use here, keep the functional contrast: \textbf{\sk{तव}} addresses one person; \textbf{\sk{युष्माकम्}} addresses several.

Source: William Dwight Whitney, \emph{Sanskrit Grammar} (1889), §§493 and 516.
\end{cousinweb}

\subsection*{Your turn}
Say these aloud:

\begin{itemize}
\raggedright
  \item several speakers say “our house” — \emph{asmākam gṛham}
  \item speaking to several people, “your house” — \emph{yuṣmākam gṛham}
  \item one addressee — \emph{tava}; several addressees — \emph{yuṣmākam}
\end{itemize}

\begin{tcolorbox}[breakable,colback=teal!4,colframe=teal!35!black,title={Before you move on}]
You speak to several people about their house. Which first word means “your, of you all”? (\textbf{\sk{युष्माकम्}}, \emph{yuṣmākam}.)
\end{tcolorbox}
\section[Practice]{One possessor or several?}

\label{lesson:SA-C259-plural-possessors-practice}



\begin{quote}
Give the three forms you can now retrieve directly: “my”, “our”, and “your, of you all”. (\emph{Mama; asmākam; yuṣmākam.})
\end{quote}

\begin{grammarlens}[title={Grammar Lens: the choice}]
Choose by counting the possessors or addressees:

\begin{itemize}
\raggedright
  \item one possessor speaking: \textbf{\sk{मम}} (\emph{mama});
  \item several possessors speaking: \textbf{\sk{अस्माकम्}} (\emph{asmākam});
  \item one familiar addressee: \textbf{\sk{तव}} (\emph{tava});
  \item several addressees: \textbf{\sk{युष्माकम्}} (\emph{yuṣmākam}).
\end{itemize}

All four answer “whose?” In this paradigm they are fixed genitive forms, \textbf{not possessive adjectives} that you inflect to agree with the thing possessed. Whitney §493 records the older adjectival history of the two plural forms; that history does not change the learner's choice here.
\end{grammarlens}

\subsection*{Your turn — choose before you speak}
Say these aloud:

\begin{itemize}
\raggedright
  \item one speaker, one family: “my family” — \emph{mama kuṭumbam}
  \item several speakers, one family: “our family” — \emph{asmākam kuṭumbam}
  \item to one familiar person: “your house” — \emph{tava gṛham}
  \item to several people: “your house” — \emph{yuṣmākam gṛham}
  \item the rule — choose by whether the possessor or addressee is one or several
\end{itemize}

\begin{tcolorbox}[breakable,colback=teal!4,colframe=teal!35!black,title={Before you move on}]
Without looking back, give “our family” and “your house” when you address several people. Then say why the first word changes. (\emph{Asmākam kuṭumbam; yuṣmākam gṛham; the possessor or addressee is several, not one.})
\end{tcolorbox}
